\chapter{Appendix}
\label{chap:appendix}

\section{Metric Command Reference}
\label{sec:appendix_commands}

Table~\ref{tab:commands} lists the subcommands of the Rust implementation, so
every result in Chapter~\ref{chap:experiments} can be reproduced. Each metric
writes \texttt{results/<label>/<metric>.json} and a matching \texttt{.csv}.

\begin{table}[htbp]
  \centering
  \caption[\entrydesc{Metric command reference}{The rust\_metrics subcommands and the module behind each}]{Subcommands of \texttt{rust\_metrics}, with
    the module implementing each. Metrics 9, 10, 12 and 13 need a second
    endpoint in \texttt{QLEVER\_ENDPOINT\_B}.}
  \label{tab:commands}
  \footnotesize
  \begin{tabular}{rlll}
    \toprule
    \# & Command & Module & Level \\
    \midrule
    1  & \texttt{population}       & \texttt{class\_population.rs}       & class \\
    2  & \texttt{entf}             & \texttt{property\_entropy.rs}       & class \\
    3  & \texttt{entetimp}         & \texttt{ent\_etimp.rs}              & class \\
    4  & \texttt{classentropy}     & \texttt{class\_entropy.rs}          & class \\
    5  & \texttt{inforank}         & \texttt{entity\_informativeness.rs} & entity \\
    6  & \texttt{diversity}        & \texttt{object\_diversity.rs}       & entity \\
    7  & \texttt{entropy-pagerank} & \texttt{entropy\_pagerank.rs}       & global \\
    8  & \texttt{shape}            & \texttt{graph\_shape.rs}            & global \\
    9  & \texttt{churn}            & \texttt{class\_churn.rs}            & global, 2 endpoints \\
    10 & \texttt{diff}             & \texttt{triple\_diff.rs}            & global, 2 endpoints \\
    11 & \texttt{trajectories}     & \texttt{trajectories.rs}            & global, stored results \\
    12 & \texttt{vocab}            & \texttt{vocab\_evolution.rs}        & global, 2 endpoints \\
    13 & \texttt{crosskg}          & \texttt{cross\_kg.rs}               & global, 2 endpoints \\
    \bottomrule
  \end{tabular}
\end{table}

\section{Environment Contract}
\label{sec:appendix_env}

Three environment variables configure every run.
\texttt{QLEVER\_ENDPOINT} names the graph to measure.
\texttt{QLEVER\_ENDPOINT\_B} names the second graph for the comparison metrics.
\texttt{QLEVER\_LABEL} names the snapshot and therefore the output directory
\texttt{results/<label>/}, which is what metric~11 reads back when it assembles a
trajectory across three or more releases.

A representative invocation, reproducing the cross-graph results of
Section~\ref{sec:exp5}. The first argument restricts the comparison to one
class; passing \texttt{-} instead matches every class name common to both
graphs, and the second argument sets how many classes are scanned per side:

\begin{verbatim}
export QLEVER_ENDPOINT=http://localhost:9006     # YAGO 4.5.0.2
export QLEVER_ENDPOINT_B=http://localhost:9010   # DBpedia 2025-12-01
export QLEVER_LABEL=yago-4.5.0.2-vs-dbpedia-2025
cargo run --release crosskg - 150
\end{verbatim}

The YAGO~4 runs add one variable, which drops matched classes by local name and
records them under \texttt{skipped} in the output:

\begin{verbatim}
export QLEVER_ENDPOINT=http://localhost:9005     # YAGO 4
export QLEVER_CROSSKG_SKIP=thing,creativework
\end{verbatim}

\section{Reproducing the Figures}
\label{sec:appendix_figures}

Every figure and table in Chapter~\ref{chap:experiments} is generated from CSV
files in \texttt{writing/data/}, exported directly from the JSON results. No
value was transcribed by hand, so re-running a metric and re-exporting is enough
to update the document.
